% Data chapter, 2 to 3 pages: the three complementary data sets. Section files in numeric order:
% 01 baseline CT volumes, 02 visual analysis (baseline and follow-up), 03 Medis centerlines and contours.
% Numbers from CGPS_data/report.md, CGPS_data/plaque_location.md and CGPS_data/report_qc.md.
% Red notes need \usepackage{xcolor}. Remove before submission.
\providecommand{\clarify}[1]{\textcolor{red}{\textbf{Clarification needed:} #1}}

\chapter{Data}
\label{ch:data}

Three data sets from the Herlev--{\O}sterbro cohort were used in this thesis, each serving a complementary role.
The baseline coronary computed tomography angiography (CCTA) volumes were the input to the in-house tracer, whose centerlines and lumen were used to compute the geometric features; the visual analysis at baseline and follow-up defined the study cohort and the outcome; and the Medis centerlines and contours served as the reference for the tracer.
A summary of the three data sets is provided in Table~\ref{tab:data-overview}.

\begin{table}[htbp]
  \centering
  \caption{The three data sets, each linked to the others by the participant's study identifier.}
  \label{tab:data-overview}
  \small
  \begin{tabular}{p{0.16\linewidth}p{0.3\linewidth}p{0.12\linewidth}rp{0.18\linewidth}}
    \toprule
    Data set & Content & Visit & Scans & Used for \\
    \midrule
    CT volumes & non-contrast and CCTA volumes & baseline & [n] & in-house tracer and geometric features \\
    Visual analysis & stenosis grade and plaque type per segment; SIS and SSS & baseline and follow-up & 2,116 & cohort and outcome \\
    Medis centerlines and contours & centerline trees labelled by segment; lumen and plaque contours of the main vessels with plaque & [visit] & 744 / [n] & reference for tracer centerlines and lumen \\
    \bottomrule
  \end{tabular}
\end{table}
